\documentclass[11pt,a5paper,landscape]{article}

\usepackage{main}
\usepackage{exercices}
\usepackage{tcolorbox}

\renewcommand{\typedocument}{Exercices}
\renewcommand{\classe}{3e}

\begin{document}

\begin{multicols*}{2}

\exercice{Schémas de chaîne d'énergie}

Dessiner les schémas de chaîne d'énergie des systèmes suivants :
\begin{enumerate}[label=\alph*)]
    \item Une pile branchée au moteur d'un drone. Les deux appareils chauffent pendant le fonctionnement.
    \item Une centrale thermique. Pour fonctionner, le charbon est brûlé dans la chaudière pour chauffer de la vapeur. La vapeur fait tourner un système turbine/alternateur pour produire de l'électricité.
\end{enumerate}

\exercice{Energies cinétique et potentielle}

\question{Calculer l'énergie cinétique des objets suivants :}
\begin{enumerate}[label=\alph*)]
    \item Un camion de masse \SI{5000}{kg} roulant à \SI{36}{km/h}.
    \item Une balle de masse \SI{30}{g} tirée à \SI{1000}{m/s}.
\end{enumerate}
\question{Calculer l'énergie potentielle de pesanteur d'une pomme de masse \SI{150}{g} située à une hauteur de \SI{2}{m} au-dessus du sol sur la Terre.}

\exercice{Saut olympique}

Une plongeuse de masse \SI{60}{kg} est immobile en haut d'un plongeoir de \SI{10}{m} de haut.

\begin{enumerate}[label=\alph*)]
    \item Calculer l'énergie potentielle de la plongeuse.
    \item Quelle est son énergie cinétique avant de sauter ? En déduire son énergie mécanique.
    \item La plongeuse saute du plongeoir. En considérant qu'il n'y a aucune perte d'énergie, quelle est son énergie mécanique juste avant de toucher l'eau ?
    \item Calculer l'énergie potentielle de pesanteur lorsqu'elle touche l'eau (hauteur = 0 m). En déduire son énergie cinétique à ce moment-là, puis sa vitesse.
\end{enumerate}

\exercice{Marathon}

Lors d'un effort physique, une transformation chimique a lieu au niveau des muscles : l'énergie chimique est convertie en énergie mécanique et en chaleur.

\question{Représenter le diagramme énergétique du muscle.}

Durant une heure de course à pied, la dépense énergétique moyenne de l’athlète est d’environ 30 kJ par kg de masse corporelle.
Une athlète de 65 kg court pendant 30 min. Pour couvrir ses besoins énergétiques, l’athlète consomme une boisson énergétique : un verre apporte une énergie d’environ 335 kJ.

\question{Déterminer le nombre de verres de boisson énergisante nécessaires pour couvrir la dépense énergétique.}

\end{multicols*}
\end{document}